\chapter{The cavity interpolation and its tree laws}\label{sec:c7}

This section and the next two prove \cref{prop:c7-closure}. Suppose that
at a minimizing profile, the scalar response operator $J^{\scal}_{a,d}$
of \cref{sec:scalar-stability} satisfies
$J^{\scal}_{a,d}\preceq(1+\eps_N)I$ with $\eps_N\le\rho/4$. The
proposition states that the self-overlap variance
$V=\int\Var(S_q)\,\nu(\dd q)$ then obeys
$(3\rho/4)V\le C_{\rm cav}[N^{-1}+A+J+H_3]$, where $C_{\rm cav}$ depends
only on $\beta$. In \cref{cor:c7-defect}, we deduce from it the bound on
the defect $\calD_N(\rho)$ that is used in \cref{sec:b7}.

The argument is a cavity interpolation on one site. We single out the
last spin and move the variance of its couplings to the other $n=N-1$
spins into its auxiliary field. At the end of the interpolation, the last
spin is independent of the others. By exchangeability, $V$ equals $n/N$
times the integrated covariance between the squared conditional
magnetization of the last spin and the self-overlap of the remaining
spins, up to an error of order $N^{-1}$. This covariance vanishes at the
independent endpoint, and its value for the full system is its derivative
there plus a Taylor remainder. In
\cref{sec:c7-endpoint}, we show that this derivative is a multiple of the
scalar response operator applied to the covariance of the self-overlap of
the remaining spins, plus a nonpositive term and explicit errors
(\cref{lem:linear}). We also bound the errors by $A$ and $J$ there
(\cref{lem:cl-response}). In \cref{sec:c7-closure}, we compare the
variance of this self-overlap at the independent endpoint with $V$, with
coefficient one (\cref{lem:cl-star}). We then express $V$ through the
expansion (\cref{lem:cav-trace}), bound the cubic moment that controls
the remainder by $A$, $J$, and $H_3$ (\cref{lem:c7-thermal}), and close the
resulting inequality.

In this section, we set up the interpolation and the calculus that the
next two sections use. In \cref{subsec:c7-row}, we define the cavity law
$\Pp_u$, $0\le u\le1$, identify its two endpoints, and record a variance
bound \eqref{eq:cl-BL} for the remaining spins. In \cref{subsec:c7-tree},
we represent products of conditional expectations by laws of finitely
many descendants on marked trees. We show that the solution of the Parisi
equation depends smoothly on the incident row and on the variance of the
last field (\cref{lem:tree-regularity}). We compute the density of the
tree laws (\cref{lem:tree-density}) and the responses
\eqref{eq:nodeTheta} and \eqref{eq:nodeb} to these two parameters
(\cref{lem:tree-responses}). From them, we derive an exact formula for the
derivative of a tree moment in $u$ (\cref{lem:streaming}) and a
comparison of tree moments at different values of $u$
(\cref{cor:comparison}). In \cref{subsec:c7-three}, we bound integrated
products of three overlap defects by a cubic moment $\mathcal M_3(u)$
(\cref{lem:cav-threeproducts}), and deduce a bound on the Taylor
remainder (\cref{lem:remainder}).

We use the generalized law of \cref{sec:gauss}, with its variance bound
(\cref{lem:fnd-posterior}) and its heat identity
(\cref{lem:fnd-path-heat}). We also use the scalar law \eqref{eq:sc-pde}
with \cref{lem:cl-scalar}, and the differentiation and approximation
arguments of \cite[Sections~2.1, 2.3, and~2.6]{AL}. The results of this
section are used as follows. The identity \eqref{eq:nodeTheta} and the
density \eqref{eq:treedensity} with one leaf enter the proof of
\cref{lem:linear}. The bound \eqref{eq:cl-BL} and \cref{cor:comparison}
enter the proofs of \cref{lem:cl-response,lem:c7-thermal}. The bound of
\cref{lem:cav-threeproducts} enters the proof of \cref{lem:cl-star}, and
\cref{lem:remainder} enters the proofs of \cref{lem:cav-trace} and
\cref{prop:c7-closure}.

The cavity method originates in the physics literature
\cite{MezardParisiVirasoro1986,MezardParisiVirasoro1987}. We use its rigorous form on
one site, the framework of differentiation, comparison, and contraction in
\cite[Section~1.6]{TalagrandVolI2011} (see Lemmas~1.6.3--1.6.4 and
Propositions~1.6.5--1.6.6 there). An expansion on one site for
hierarchical Gaussian laws appears in \cite[Lemma~2.2]{OpenAILimitingLaw}.
There, a site is chosen by exchangeability and deleted, and its Gaussian
kernels are interpolated to those of an independent scalar model on one
site. The linear cavity kernel is half of a connected correlation of four
spins of the removed site \cite[(A8)]{OpenAILimitingLaw}, and a linear
equation with this kernel is inverted under a gap condition on the scalar
response in \cite[Proposition~5.1]{OpenAILimitingLaw}. The tree laws of
\cref{subsec:c7-tree} and their calculus are adapted from
\cite[Section~2]{OpenAIFluctuationScale}; we give the precise references
there.

\paragraph{Standing assumptions.}
Throughout \cref{part:cavity}, we fix $N\ge2$ and write $n=N-1$. We also
fix $0<\rho\le1$, $t=1-\rho$, $c=\beta^2\rho$, and
$0<s=\lambda/\beta\le1/2$. In this section, the profile
$\mu=s\delta_0+\nu$ is deterministic and otherwise arbitrary. The measure
$\nu$ may have atoms anywhere in $[0,1]$, including $0$ and $1$. No result
of this section uses a minimizing property. This matters because at the
independent endpoint, the remaining $n$ spins form a system whose law is
built from the profile of the full system, and this profile need not
minimize any functional of the $n$-spin system. Constants $C$ are absolute, and
constants $C_\beta$ depend only on $\beta$. Neither depends on $N$,
$\rho$, $s$, the profile, or the positions and masses of its atoms. Note
that the interpolation parameter $u$ is not the solution $u$ of
\eqref{eq:sc-pde}; in this part of the paper, we write $U^0$ for the
latter.

\section{Removing the incident row}\label{subsec:c7-row}

We write $\sigma=(\tau,\epsilon)\in\{-1,1\}^n\times\{-1,1\}$, so that
$\epsilon=\sigma_N$ is the last spin. The couplings $g_{iN}$, $i\le n$,
between $\epsilon$ and $\tau$ form the \emph{incident row}. Let
$u\in[0,1]$ be the interpolation parameter; $t$ stays fixed. We set
\begin{equation}\label{eq:cavityparameters}
 \alpha=\beta^2t\,\frac nN,\qquad
 d=c+\alpha=\beta^2\Bigl(1-\frac tN\Bigr),\qquad
 b(u)=d-\alpha u,\qquad 0\le u\le1.
\end{equation}
The second identity holds because
$c+\alpha=\beta^2(\rho+t-t/N)=\beta^2(1-t/N)$. Since $0\le t<1$ and
$N\ge2$, we have $d\in[\beta^2/2,\beta^2]$. Further, $b(u)\ge b(1)=c>0$
for every $u$. For $x\in\R^n$ and $y\in\R$, let
\begin{gather*}
 \mathcal E(\sigma;x,y)=\mathcal E^-(\tau)
 +\sqrt{\frac{\alpha u}{n}}\;\epsilon\sum_{i=1}^ng_{iN}\tau_i
 +x\cdot\tau+y\epsilon,
 \\
 \mathcal E^-(\tau)=\frac{\beta\sqrt t}{\sqrt N}
       \sum_{1\le i<j\le n}g_{ij}\tau_i\tau_j,
\end{gather*}
where the couplings are independent standard Gaussian variables under the
reference law. Let $U(q,x,y)$ solve
\begin{equation}\label{eq:cavityPDE}
 \begin{gathered}
 \partial_q U=-\frac12\Bigl\{
 c\bigl(\Delta_xU+a(q)|\nabla_xU|^2\bigr)
 +b(u)\bigl(\partial_{yy}U+a(q)(\partial_yU)^2\bigr)\Bigr\},
 \\ U(1,x,y)=\log\sum_\sigma e^{\mathcal E(\sigma;x,y)}.
 \end{gathered}
\end{equation}
At $u=1$, each coupling of the row has the variance $\alpha/n=\beta^2t/N$
of the full system, and the last field has the variance coefficient
$b(1)=c$ of the full system. As $u$ decreases to $0$, the variance
$\alpha u/n$ of each coupling in the row is removed, and the coefficient
of the last field increases to $b(0)=d=c+\alpha$.

The completion of \cref{sec:fnd} adds to the terminal exponent the
Gaussian term $G=\beta\sqrt t\sum_{i\le N}z_i/\sqrt{2N}$, which does not
depend on the spins. It changes $U$ into $U+G$ and multiplies the tilt
below by $e^{sG}$. Since $G$ is independent of the couplings $g_{ij}$,
this factor becomes, after normalization, an independent tilt of the
variables $z_i$, and these variables do not enter the field or the
spins. We omit it.

We construct the law $\Pp_u$, with expectation $\E_u$, as follows. We
tilt the reference law of the couplings by $e^{sU(0,0,0)}/Z_s(u)$,
where $Z_s(u)=\E_{\rm ref}e^{sU(0,0,0)}$. Given the couplings, the field
path $(X_q,Y_q)\in\R^n\times\R$ solves
\[
 \begin{aligned}
 \dd X_q&=c\,a(q)\nabla_xU(q,X_q,Y_q)\dd q+\sqrt c\,\dd B_q,
 &X_0&=0,\\
 \dd Y_q&=b(u)\,a(q)\partial_yU(q,X_q,Y_q)\dd q+\sqrt{b(u)}\,\dd B'_q,
 &Y_0&=0,
 \end{aligned}
\]
where $(B,B')$ is a standard Brownian motion in $\R^{n+1}$. The diffusion
matrix is $\operatorname{diag}(cI_n,b(u))$, and the drift is $a(q)$
times this matrix applied to $\nabla U(q,X_q,Y_q)$. Given the path,
we draw $\sigma$ with probabilities proportional to
$e^{\mathcal E(\sigma;X_1,Y_1)}$. Let $\calG_q$ be the $\sigma$-field
generated by the couplings and the path up to level $q$.

\paragraph{The two endpoints.}
At $u=1$, the law $\Pp_1$ is the tilted law of \cref{subsec:fnd-law} for
the full system, with the completion omitted as above. At $u=0$, the row
is absent, and the terminal value separates as
$U(1,x,y)=U^-(1,x)+\log(2\cosh y)$, where
$U^-(1,x)=\log\sum_\tau e^{\mathcal E^-(\tau)+x\cdot\tau}$. Let $U^-$ be
the solution for the $n$-spin system with Hamiltonian $\mathcal E^-$,
field coefficient $c$, and profile $\mu$, and let $U^0$ solve
\eqref{eq:sc-pde} with coefficient $d$ and profile $\mu$. We claim that
$U=U^-+U^0$ at $u=0$. For a step profile, each Hopf--Cole step of
\eqref{eq:cavityPDE} factors, because the increments of $x$ and of $y$
are independent and $b(0)=d$. Applying the estimate \cite[(2.19)]{AL} to
each of the three equations, we extend the identity to every profile
through the approximation of \cite[Section~2.6]{AL}. The outer normalizer
factors as $Z_s(0)=e^{sU^0(0,0)}\E_{\rm ref}e^{sU^-(0,0)}$, and the
deterministic factor $e^{sU^0(0,0)}$ cancels from the tilt. The drift of
$X$ involves only $\nabla_xU^-$, and that of $Y$ only $\partial_yU^0$.
The two Brownian motions are independent, and the terminal Gibbs
probability is the product of
$e^{\mathcal E^-(\tau)+X_1\cdot\tau-U^-(1,X_1)}$ and
$e^{\epsilon Y_1}/(2\cosh Y_1)$. We conclude that $\Pp_0$ is the product
of the tilted law of the $n$-spin system, with expectation $\E^-$, and
the scalar law of \cref{sec:scalar-stability} with coefficient $d$ and
profile $\mu$, with expectation $\E^0$. Since $d\in[\beta^2/2,\beta^2]$,
we may use the scalar quantities $m_p$, $h_p$, $f_p=m_p^2$, $\bar f_p=f_p-\E^0f_p$,
$g_r^\ell$, and $\Gamma(p)=\Gamma_{a,d}(p)=\E^0m_p^2$ of
\cref{sec:scalar-stability} for this scalar law, together
with the bounds of \cref{lem:cl-scalar}.

\paragraph{The remaining spins.}
For the first $n$ spins, we define
\[
 M_q=\E_u[\tau\mid\calG_q],\qquad A_q^-=\frac{|M_q|^2}{n},\qquad
 C_q^-=\Cov_u(\tau\mid\calG_q),\qquad
 H_q^-=\nabla_x^2U(q,X_q,Y_q).
\]
We let $m_N(q)=\E_u[\epsilon\mid\calG_q]$ be the conditional
magnetization of the last spin, and we set
\begin{equation}\label{eq:cav-rest-errors}
 a_q^-(u)=\frac{\E_u\Tr\bigl((C_q^-)^2\bigr)}{n^2},\qquad
 j_q^-(u)=\frac{\E_u[M_q^TC_q^-M_q]}{n^2}.
\end{equation}
The \emph{rest self-overlap} $A_q^-$ is a random variable. It should not
be confused with the trace term $A$ of the full system.

We note that $\Pp_u$ is a generalized law in the sense of \cref{sec:gauss}.
It has $N$ coordinates, the diffusion matrix
$\Lambda=\operatorname{diag}(cI_n,b(u))$, and independent Gaussian
couplings with deterministic coefficients: $\beta\sqrt{t/N}$ for the
pairs $i<j\le n$, $(\alpha u/n)^{1/2}$ for the pairs $(i,N)$, and $0$ on
the diagonal. The matrix $\Lambda$ is positive definite because $c>0$
and $b(u)\ge c$. The coefficient $(\alpha u/n)^{1/2}$ vanishes at $u=0$,
which the generalized law allows. We use \cref{lem:fnd-posterior} for this
law, with the first $n$ coordinates selected. It implies that
$M_q=\nabla_xU(q,X_q,Y_q)$ and $m_N(q)=\partial_yU(q,X_q,Y_q)$. Since
principal submatrices preserve the order $a(q)C_q\preceq H_q\preceq C_q$,
it also implies that $a(q)C_q^-\preceq H_q^-\preceq C_q^-$, and in
particular that $0\preceq H_q^-\preceq C_q^-$. Taking $k=\psi$ in the two
inequalities in \eqref{eq:fnd-posterior}, we find that for every
$u\in[0,1]$ and every deterministic $\psi\in L^2(\nu)$,
\begin{equation}\label{eq:cl-BL}
 \Var_u\Bigl(\int\psi(q)A_q^-\,\nu(\dd q)\Bigr)
 \le4\int\psi(q)^2j_q^-(u)\,\nu(\dd q).
\end{equation}

\section{Laws of finitely many descendants}\label{subsec:c7-tree}

Derivatives in $u$ produce products of conditional expectations at
several levels. We represent them by additional terminal spins, called
leaves, under a joint law in which deleting a leaf does not change the
law of the others. The hierarchical background of such laws is Ruelle's
probability cascades \cite{Ruelle1987}; see
\cite[Chapter~2]{Panchenko2013Book}. An earlier hierarchical Gaussian
change of measure appears in
\cite[Theorem~4.3, (4.25)--(4.26)]{BolthausenSznitman1998}. Normalized
finite replica trees of the present kind were developed in
\cite[Section~2]{OpenAIFluctuationScale}, together with the
differentiation of their moments by attaching new leaves and the
resulting Taylor and comparison bounds. They are also used in
\cite[Sections~2.1--2.3]{OpenAILimitingLaw}. In
\cite{OpenAIFluctuationScale}, the covariance differentiation formula is
Proposition~2.5, the attachment rules and their uniform variation bounds
are Lemma~2.6, the bounds on the levels at which new pairs split are
Lemma~2.8, and the Taylor and comparison bounds are Lemma~2.9. We derive
below the versions that we need for the cavity law, with its root tilt
and an arbitrary profile, including its atoms.

\subsection{Marked trees}

\begin{definition}[Marked trees and tree laws]\label{def:tree-marked}
A \emph{marked tree} $\mathcal T$ with $k\ge1$ leaves is a finite rooted
tree whose vertices $v$ carry deterministic levels $q_v\in[0,1]$, with
the following properties. The root has level $0$, and levels increase
strictly along every edge. Every vertex at a level in $(0,1)$ has at
least two outgoing edges. The vertices at level $1$ have no outgoing
edges; we call them terminal points. Finally, $k$ labeled leaves
$\mathsf a=1,\ldots,k$ are attached to the terminal points, and every
terminal point carries at least one leaf.

Given $u\in[0,1]$, the \emph{tree law} $\Pp_{\mathcal T,u}$, with
expectation $\E_{\mathcal T,u}$, is constructed as follows. The couplings
are drawn once from their tilted law under $\Pp_u$. Given the couplings,
the field starts at the origin at the root and follows the diffusion of
\cref{subsec:c7-row} along each edge. At each vertex, it continues along
the outgoing edges in conditionally independent ways. Given the couplings
and the field, the leaves carry conditionally independent spins
$\sigma^{\mathsf a}=(\tau^{\mathsf a},\epsilon_{\mathsf a})$. The spin
$\sigma^{\mathsf a}$ is drawn with probabilities proportional to
$e^{\mathcal E(\sigma;X^{\mathsf a}_1,Y^{\mathsf a}_1)}$, where
$(X^{\mathsf a}_1,Y^{\mathsf a}_1)$ is the field at the terminal point of
$\mathsf a$. We write $\E_{k,q,u}$ for the tree law of $k$ leaves that
all split at the level $q$.
\end{definition}

The field is defined at every point $v$ of $\mathcal T$, that is, at every
vertex and at every point of an edge. We write $(X_v,Y_v)$ for it and
$q_v$ for the level of $v$. Two points on different branches at the same
level are different points, with different fields. If $k\ge2$ and
$0<q<1$, the tree of $k$ leaves that split at $q$ has one edge from the
root to a vertex at level $q$, followed by $k$ edges to $k$ terminal
points. If $k\ge2$ and $q=0$, the root has $k$ outgoing edges, and if
$k\ge2$ and $q=1$, the $k$ leaves share one terminal point. For $k=1$,
the tree is the single path from level $0$ to level $1$, independently of
$q$. After deleting leaves, we remove the branches that carry no
remaining leaf and suppress the interior vertices with one outgoing edge.

Given the couplings and the field at a vertex, the continuations along
different outgoing edges are conditionally independent, and their laws are
normalized. Integrating out a leaf yields the tree law of the
remaining leaves. We call this the \emph{deletion property}. In
particular, every leaf has law $\Pp_u$. We represent products of conditional
expectations given $\calG_q$ by leaves that split at $q$.
For leaves $\mathsf a\ne\mathsf b$, let $q_{\mathsf{ab}}$ be the level of
their last common ancestor, so that $q_{\mathsf{ab}}=1$ if they share a
terminal point. We define their \emph{rest overlap}, their \emph{overlap
defect}, and the product of the overlap defect with their last spins by
\begin{equation}\label{eq:defectdefinition}
 R^-_{\mathsf{ab}}=\frac{\tau^{\mathsf a}\cdot\tau^{\mathsf b}}{n},\qquad
 \Delta_{\mathsf{ab}}=R^-_{\mathsf{ab}}-q_{\mathsf{ab}},\qquad
 D_{\mathsf{ab}}=\epsilon_{\mathsf a}\epsilon_{\mathsf b}
                          \Delta_{\mathsf{ab}}.
\end{equation}
Since $|R^-_{\mathsf{ab}}|\le1$, we have
$|D_{\mathsf{ab}}|=|\Delta_{\mathsf{ab}}|\le2$. For instance, the four
spins of $\E_{4,p,u}$ are conditionally independent given $\calG_p$, each
with the conditional law of $\sigma$ given $\calG_p$. Therefore
$\E_u[m_N(p)^2A_p^-]=\E_{4,p,u}[\epsilon_1\epsilon_2R^-_{34}]$.

\subsection{Regularity in the row and in the last-field variance}

To differentiate tree moments, we temporarily replace the row by a
deterministic coupling $\epsilon\sum_i\theta_i\tau_i$ with
$\theta\in\R^n$, and the coefficient $b(u)$ by a number $b>0$. We write
$\Lambda_b=\operatorname{diag}(cI_n,b)$. We keep the other couplings $g_{ij}$,
$i<j\le n$, fixed. Under $\Pp_u$ and under the tree laws,
conditionally on the couplings, the field and the spins have the law of
this fixed system with $\theta=(\alpha u/n)^{1/2}(g_{1N},\ldots,g_{nN})$
and $b=b(u)$.

\begin{lemma}[Regularity in the parameters]\label{lem:tree-regularity}
Let $\mu$ be a deterministic profile, and let $U=U(q,x,y)$ solve
\eqref{eq:cavityPDE} with the row replaced by
$\epsilon\sum_i\theta_i\tau_i$ and $b(u)$ replaced by $b>0$. Then $U$ has
derivatives of every order in $(\theta,x,y)$. Each derivative of positive
order is bounded by a constant depending only on $N$ and its order. This
bound is uniform in the couplings, $\theta$, $b$, $(q,x,y)$, and the
profile, and each first derivative is bounded by one in absolute value.
The function $U$ and all these derivatives are jointly continuous in
$(b,q,\theta,x,y)\in(0,\infty)\times[0,1]\times\R^n\times\R^n\times\R$.
Finally, writing $\nabla$ for the gradient in $(x,y)$, we have for
$0\le r\le h\le1$ that
\begin{equation}\label{eq:tree-integrated-pde}
 U(r,x,y)-U(h,x,y)=\frac12\int_r^h\bigl[\Tr(\Lambda_b\nabla^2U)
   +a(q)\,\nabla U^T\Lambda_b\nabla U\bigr](q,x,y)\dd q.
\end{equation}
\end{lemma}

\begin{proof}
The proof has three steps. We first prove the bounds for step profiles,
using the density of the future of the field. We then pass to arbitrary
profiles with the approximation of \cite[Section~2.6]{AL}, and finally
we prove the continuity in $q$ and $b$.

\paragraph{Step 1: step profiles.}
Let $\mu$ be a step profile. The coefficients of \eqref{eq:cavityPDE} do
not depend on $(x,y)$. Translating the starting point of the field by
$(x',y')$ has the same effect as adding
$x'\cdot\tau+y'\epsilon$ to the terminal exponent. It follows that every
derivative in $(\theta,x,y)$ is a derivative at $z=0$ with respect to
parameters $z_j$ of the terminal exponent
\[
 \mathcal E_z(\sigma;x,y)=\mathcal E(\sigma;x,y)+\sum_jz_j\eta_j(\sigma),
\]
where the features $\eta_j$ range over $\tau_i\epsilon$, $\tau_i$
($i\le n$), and $\epsilon$. Note that $|\eta_j|\le1$. Let $U_z$ be the
corresponding solution. We fix a level $q$, refine the grid of the profile
so that it contains $q$, and fix the field $(x,y)$ at level $q$. On a
grid interval $[q_\ell,q_{\ell+1})$ with $a=a_\ell>0$, the Hopf--Cole form
of the equation is $e^{a_\ell U_z(q_\ell)}=\E e^{a_\ell U_z(q_{\ell+1})}$,
where the expectation is over a Gaussian increment of covariance
$(q_{\ell+1}-q_\ell)\Lambda_b$. The transition of the diffusion is this
Gaussian transition tilted by $e^{a_\ell(U_z(q_{\ell+1})-U_z(q_\ell))}$.
If $a_\ell=0$, it is the Gaussian transition. We multiply these factors
and the Gibbs probability $e^{\mathcal E_z(\sigma)-U_z(1)}$, as in
\cite[Section~2.1, (2.1)]{AL}. The coefficient of $U_z$ at a grid point
$q_\ell>q$ is $a_{\ell-1}-a_\ell=-\nu(\{q_\ell\})$, and at level $1$ it is
$a(1-)-1=-\nu(\{1\})$. We find that given the field at level $q$, the
future increments of the field and the spin have the density
\begin{equation}\label{eq:tree-future-density}
 L_z=\exp\Bigl\{\mathcal E_z(\sigma;Z_1)-a(q)U_z(q,Z_q)
     -\int_{(q,1]}U_z(r,Z_r)\,\nu(\dd r)\Bigr\}
\end{equation}
with respect to driftless Gaussian increments of covariance $\Lambda_b$
and counting measure on $\sigma$. Here $Z_r$ is the field at level $r$.

We write $\E_z$ for the expectation under this density and
$\ell_z=\log L_z$. On a grid interval with $a_\ell>0$, we differentiate
$e^{a_\ell U_z(q_\ell)}=\E e^{a_\ell U_z(q_{\ell+1})}$. This shows that
$\partial_{z_j}U_z(q_\ell)$ is the expectation of
$\partial_{z_j}U_z(q_{\ell+1})$ under the transition of the diffusion,
and the same holds if $a_\ell=0$. At level $1$, $\partial_{z_j}U_z(1)$ is
the Gibbs mean of $\eta_j$. Iterating, we obtain
\begin{equation}\label{eq:tree-first-derivative}
 \partial_{z_j}U_z(q)=\E_z\eta_j,\qquad\text{so that}\qquad
 |\partial_{z_j}U_z|\le1 .
\end{equation}
Since $a(q)+\nu((q,1])=1$, this implies that $|\partial_{z_j}\ell_z|\le2$.
Further, if all derivatives of $U_z$ of some order $m'\ge2$ are bounded
by $C_{m'}$, then all derivatives of $\ell_z$ of order $m'$ are bounded
by $C_{m'}$. Since the reference law in \eqref{eq:tree-future-density} does not
depend on $z$, we may differentiate \eqref{eq:tree-first-derivative}
under the expectation. For directions $j,j_1,\ldots,j_m$, we have
\[
 \partial_{z_{j_1}}\cdots\partial_{z_{j_m}}\partial_{z_j}U_z(q)
 =\E_z\Bigl[\eta_j\sum_\pi\prod_{B\in\pi}\partial_B\ell_z\Bigr],
\]
where $\pi$ runs over the partitions of $\{1,\ldots,m\}$ and $\partial_B$
is the derivative in the directions $j_i$, $i\in B$. By induction on $m$,
we find that every derivative of order $m+1$ in the coordinates $z_j$ is
bounded by a constant depending only on $m$. There are $2n+1$
coordinates, and this proves the bounds of the lemma for step profiles.

\paragraph{Step 2: arbitrary profiles.}
As in \cite[Section~2.6.2]{AL}, we approximate $\mu$ by step profiles
$\mu_m=s\delta_0+\nu_m$ that converge weakly to $\mu$ and whose
cumulative functions $a_m$ converge to $a$ in $L^1(0,1)$. Let $U_m$ be
the corresponding solutions. Following the proof of \cite[(2.19)]{AL}, with $cI_N$
replaced by $\Lambda_b$ and with the first derivatives bounded by one,
we obtain
\[
 \sup|U_m-U|\le\frac{\Tr\Lambda_b}2\,\|a_m-a\|_{L^1(0,1)},
\]
where the supremum is over $q,x,y$, the couplings, and $\theta$. For a
function $f$ of one real variable and every $h>0$, we have
$\|f'\|_\infty\le2\|f\|_\infty/h+(h/2)\|f''\|_\infty$. Using this
inequality in one direction at a time for $U_m-U_{m'}$ and then for its
derivatives, together with the bounds of Step 1, we find that every derivative of $U_m$ in
$(\theta,x,y)$ converges uniformly. This implies that $U$ has
derivatives of every order in $(\theta,x,y)$, that they are the uniform
limits of those of $U_m$, and that they satisfy the same bounds. Each $U_m$
satisfies \eqref{eq:tree-integrated-pde} with $a_m$ in place of $a$. As
$a_m\to a$ in $L^1$, the identity passes to the limit.

\paragraph{Step 3: continuity.}
For a step profile, the derivatives of $U$ in $(\theta,x,y)$ satisfy, on
each grid interval, the equations obtained by differentiating
\eqref{eq:cavityPDE}. By Step 1, their right sides are bounded by
constants depending only on $N$, $c$, $b$, and the order. More precisely,
for every multi-index $\kappa$ in $(\theta,x,y)$, including $\kappa=0$,
differentiation and the product rule give
\[
 \sup\bigl|\partial_qD^\kappa U_m\bigr|\le C_{N,|\kappa|}\,(nc+b)
\]
on each grid interval, where $C_{N,|\kappa|}$ depends only on $N$ and
$|\kappa|$. The derivatives are continuous across the grid points, so the
same Lipschitz bound in $q$ holds on $[0,1]$, independently of the step
approximation, and it passes to the limit. We also find that
$U$ and its derivatives in $(\theta,x,y)$ are Lipschitz in $q$,
with constants depending only on $N$, the order, and $nc+b$. For the continuity in
$b$, let $U_b$ and $U_{b'}$ be the solutions for $b,b'>0$ and the same
profile. The difference $U_b-U_{b'}$ vanishes at $q=1$, because both solutions
have the same terminal value. Subtracting their equations, we see that it
satisfies a linear backward equation with diffusion matrix
$\Lambda_{b'}$, a bounded drift, and the source term
$\frac12(b-b')\bigl[\partial_{yy}U_b+a(\partial_yU_b)^2\bigr]$. By Step 1,
we have $|\partial_{yy}U_b|\le2$ and $|\partial_yU_b|\le1$. The backward
representation then gives $\sup|U_b-U_{b'}|\le\frac32|b-b'|$, first for
step profiles and then, by Step 2, for $\mu$. Using the interpolation
inequality of Step 2 again, we find that every derivative in
$(\theta,x,y)$ is continuous in $b$, uniformly in the other variables.
These derivatives are also Lipschitz in $(\theta,x,y)$ and in $q$, so
they are jointly continuous.
\end{proof}

\subsection{The density of a tree law}

Let $\mathcal T$ be a marked tree. For a vertex $v$, let $l_v$ be the
number of its outgoing edges if $q_v<1$, and the number of leaves at $v$
if $v$ is a terminal point. We define the measure $\omega$ on the points
of $\mathcal T$ as follows. On the interior of an edge whose endpoints
have levels $r<h$, it is the copy of the restriction of $\nu$ to $(r,h)$.
A vertex $v$ carries the atom
\[
 \omega_v=\omega(\{v\})=\nu(\{q_v\})+(l_v-1)\,a(q_v).
\]
For example, the root carries $\nu(\{0\})+(l-1)a(0)$ if it has $l$
outgoing edges, and a terminal point with $l$ leaves carries
$\nu(\{1\})+l-1$. Below, the notation $\sum_v\omega_vF(v)$ means
$\int F\dd\omega$; for an atomic profile, it is an ordinary sum. The
\emph{driftless Gaussian field} on $\mathcal T$ starts at the origin at
the root and has independent increments along the edges. On each edge,
its increments have covariance $\Lambda$ times the length of the level
interval.

\begin{lemma}[Density of a tree law]\label{lem:tree-density}
Let $\mu$ be a deterministic profile, $u\in[0,1]$, and $\mathcal T$ a
marked tree with $k$ leaves. Consider the product of the original
Gaussian law of the couplings, the driftless Gaussian field on
$\mathcal T$, and counting measure on the spins. With respect to this
reference measure, the joint law of the couplings, the field on
$\mathcal T$, and the spins of the leaves under $\Pp_{\mathcal T,u}$ has
density $Z_s(u)^{-1}\mathscr D_{\mathcal T}$, where
\begin{equation}\label{eq:treedensity}
 \mathscr D_{\mathcal T}=\exp\Bigl\{
       \sum_{\mathsf a=1}^k\mathcal E(\sigma^{\mathsf a};
                  X^{\mathsf a}_1,Y^{\mathsf a}_1)
       -\int_{\mathcal T}U(q_v,X_v,Y_v)\,\omega(\dd v)\Bigr\},
 \qquad \omega\ge0,\qquad \omega(\mathcal T)=k-s.
\end{equation}
For fixed couplings, the conditional density of the field and the spins
is $e^{-sU(0,0,0)}\mathscr D_{\mathcal T}$.
\end{lemma}

\begin{proof}
The proof has three steps: we compute the likelihood of one edge, collect
the coefficients of $U$ at the vertices, and compute the mass of
$\omega$. We fix the couplings. By \cref{lem:tree-regularity}, the first
derivatives of $U$ in $(x,y)$ are bounded by one. The drift
$a(q)\Lambda\nabla U$ is bounded, with
$a(q)^2|\Lambda^{1/2}\nabla U|^2\le nc+b(u)$, and Novikov's condition
holds.

\paragraph{Step 1: one edge.}
Let $Z$ be the driftless field on an edge with levels $r<h$, so that
$\dd Z_q=\Lambda^{1/2}\dd B_q$. By Girsanov's formula, the law of the
diffusion on this edge, started from $Z_r$, has density
\[
 \exp\Bigl\{\int_r^ha(q)\nabla U(q,Z_q)\cdot\dd Z_q
   -\frac12\int_r^ha(q)^2|\Lambda^{1/2}\nabla U(q,Z_q)|^2\dd q\Bigr\}
\]
with respect to the law of $Z$. By \cref{lem:tree-regularity}, $U$ is
smooth in $(x,y)$ with bounded derivatives and satisfies
\eqref{eq:tree-integrated-pde}, whose integrand is bounded; in
particular, $U$ is Lipschitz in $q$. We may therefore apply Itô's formula
(for instance after smoothing $U$ in $q$), and by
\eqref{eq:tree-integrated-pde}, we have
\[
 \dd U(q,Z_q)=\nabla U(q,Z_q)\cdot\dd Z_q
       -\frac{a(q)}2|\Lambda^{1/2}\nabla U(q,Z_q)|^2\dd q.
\]
In the last identity, we used that the time derivative of $U$ is given by
\eqref{eq:cavityPDE} for almost every $q$. We conclude that the exponent of the density is
$\int_r^ha(q)\dd U(q,Z_q)$. The process $U_q=U(q,Z_q)$ is continuous,
and $a$ is right-continuous with Stieltjes measure $\nu$ on $(r,h]$. We
integrate by parts. The countably many jumps of $a$ do not affect the
integral against $U_q$, and we obtain
$\int_r^ha\dd U_q=a(h)U_h-a(r)U_r-\int_{(r,h]}U_q\,\nu(\dd q)$. As
$a(h)-\nu(\{h\})=a(h-)$, the edge likelihood is
\begin{equation}\label{eq:cav-edge-likelihood}
 \exp\Bigl\{a(h-)U(h,X_h,Y_h)-a(r)U(r,X_r,Y_r)
                    -\int_{(r,h)}U(q,X_q,Y_q)\,\nu(\dd q)\Bigr\},
\end{equation}
where we wrote the field as $(X,Y)$. We write the values at the endpoints
separately, so that this formula also covers atoms of $\nu$ at $r$ and
$h$.

\paragraph{Step 2: the vertices.}
We multiply the root tilt $e^{sU(0,0,0)}/Z_s(u)$, the edge likelihoods
\eqref{eq:cav-edge-likelihood}, and the Gibbs probabilities
$e^{\mathcal E(\sigma^{\mathsf a};X^{\mathsf a}_1,Y^{\mathsf a}_1)
-U(1,X^{\mathsf a}_1,Y^{\mathsf a}_1)}$ of the leaves. The open edges
contribute $-\int U\dd\nu$ along each edge. At a vertex $v$ with
$0<q_v<1$, the incoming edge contributes $a(q_v-)U(v)$, and the $l_v$
outgoing edges contribute $-l_va(q_v)U(v)$. The coefficient of $U(v)$ is
\[
 a(q_v-)-l_va(q_v)=-\bigl[\nu(\{q_v\})+(l_v-1)a(q_v)\bigr]=-\omega_v .
\]
At the root, the tilt and the $l$ outgoing edges give
$s-la(0)=-[\nu(\{0\})+(l-1)a(0)]$, since $a(0)=s+\nu(\{0\})$. At a
terminal point with $l$ leaves, the incoming edge and the $l$ Gibbs
probabilities give $a(1-)-l=-[\nu(\{1\})+l-1]$, since $a(1)=1$. All these
coefficients are nonpositive, which shows that $\omega\ge0$, and the
product is $Z_s(u)^{-1}\mathscr D_{\mathcal T}$ with
$\mathscr D_{\mathcal T}$ as in \eqref{eq:treedensity}. The edge
likelihoods and the Gibbs probabilities are normalized. Integrating
$\mathscr D_{\mathcal T}$ over the driftless field and the spins thus
gives $e^{sU(0,0,0)}$, and this proves the claim about fixed couplings.
Averaging over the original law of the couplings, we obtain $Z_s(u)$; in
particular, the normalizer is $Z_s(u)$ for every tree.

\paragraph{Step 3: the mass.}
For a single leaf, $\omega$ is a copy of $\nu$ on the path from the root
to the terminal point, and $\omega(\mathcal T)=\nu([0,1])=1-s$. Every
marked tree is obtained from a single path by successively attaching
$l-1\ge1$ new continuations at a point at some level $q$, each running to
level $1$. Such a step adds $l-1$ leaves. It adds $(l-1)a(q)$ to the
weight of the point, which is or becomes a vertex, and $(l-1)\nu((q,1])$
on the new continuations. Since $a(q)+\nu((q,1])=1$, the mass of $\omega$
increases by $l-1$, and $\omega(\mathcal T)=1-s+(k-1)=k-s$.
\end{proof}

We also see from the density \eqref{eq:treedensity} that an \emph{extended}
tree, obtained by attaching new leaves to $\mathcal T$, carries the
weights of its own shape, and that by deletion it induces
$\Pp_{\mathcal T,u}$ on the original leaves. The same proof applies when
the row is a deterministic coupling $\epsilon\sum_i\theta_i\tau_i$ and
$b(u)$ is replaced by a number $b>0$, as in \cref{lem:tree-regularity}.
We then take the density with respect to the Gaussian law of the
couplings $g_{ij}$, $i<j\le n$, and the normalizer is the expectation of
$e^{sU(0,0,0)}$ over these couplings.

\subsection{Responses to the row and to the last-field variance}

Let the field $(X,Y)$ and the spin $\sigma$ have the law of the fixed
system of \cref{lem:tree-regularity}, and let $\calG_q$ be generated by
the field up to level $q$. We set
\[
 \pi_i(q)=\E[\tau_i\epsilon\mid\calG_q],\qquad
 m_N(q)=\E[\epsilon\mid\calG_q],
\]
and we evaluate the derivatives of $U$ at $(q,X_q,Y_q)$, writing for
instance $\partial_{\theta_i}U(q)$.

\begin{lemma}[Parameter responses]\label{lem:tree-responses}
Let $\mu$ be a deterministic profile, and let $U$, $\pi_i$, and $m_N$ be
as above, for fixed couplings, $\theta\in\R^n$, and $b>0$. Then for every
$q\in[0,1]$ and $i\le n$,
\begin{equation}\label{eq:nodeTheta}
 \partial_{\theta_i}U(q)=\pi_i(q),\qquad
 \partial_{\theta_i}^2U(q)
 =1-a(q)\pi_i(q)^2-\int_{(q,1]}\E[\pi_i(r)^2\mid\calG_q]\,\nu(\dd r),
\end{equation}
\begin{equation}\label{eq:tree-node-y}
 \partial_yU(q)=m_N(q),\qquad
 \partial_{yy}U(q)
 =1-a(q)m_N(q)^2-\int_{(q,1]}\E[m_N(r)^2\mid\calG_q]\,\nu(\dd r).
\end{equation}
Further, the derivative of $U$ in $b$ exists and, for each fixed $q$, is
jointly continuous in $(b,\theta,x,y)$. It satisfies
\begin{equation}\label{eq:nodeb}
 \partial_bU(q)=\frac12\Bigl[(1-q)-\int_{(q,1]}(r-q)
                   \E[m_N(r)^2\mid\calG_q]\,\nu(\dd r)\Bigr].
\end{equation}
In particular, $0\le\partial_bU(q)\le(1-q)/2$.
\end{lemma}

For a single row, \eqref{eq:nodeTheta} and \eqref{eq:tree-node-y} are
the conditional form of \cite[(2.10)]{AL}. The identity \eqref{eq:nodeb}
is the conditional, one-coordinate form of the field-variance computation
following \cite[(2.11)]{AL}. Under $\Pp_u$ and the tree laws, the
identities hold with $\calG_q$ generated by the couplings and the field,
with $\theta=(\alpha u/n)^{1/2}(g_{1N},\ldots,g_{nN})$ and $b=b(u)$.

\begin{proof}
The proof has four steps. We prove \eqref{eq:nodeTheta} and
\eqref{eq:tree-node-y} for step profiles with the density
\eqref{eq:tree-future-density}, and we pass to arbitrary profiles with
the exact density of \cref{lem:tree-density}. We then treat the
derivative in $b$ in the same two stages.

\paragraph{Step 1: step profiles.}
Let $\mu$ be a step profile, and let $\eta$ be one of the features
$\tau_i\epsilon$ and $\epsilon$, with parameter $z$. By
\eqref{eq:tree-first-derivative}, we have
$\partial_zU(r)=\E[\eta\mid\calG_r]$ for every $r$, and we write
$\varpi(r)$ for this conditional expectation. Differentiating
\eqref{eq:tree-first-derivative} once more under the density
\eqref{eq:tree-future-density}, we find
\[
 \partial_z^2U(q)=\E\bigl[\eta\,\partial_z\ell_z\bigm|\calG_q\bigr]
 =\E[\eta^2\mid\calG_q]-a(q)\varpi(q)^2
   -\int_{(q,1]}\E\bigl[\eta\,\partial_zU(r)\bigm|\calG_q\bigr]\nu(\dd r).
\]
For $r\ge q$, by the tower property, we have
$\E[\eta\,\partial_zU(r)\mid\calG_q]=\E[\varpi(r)\E[\eta\mid\calG_r]\mid\calG_q]
=\E[\varpi(r)^2\mid\calG_q]$. Since $\eta^2=1$, this proves
\eqref{eq:nodeTheta} for $\eta=\tau_i\epsilon$ and
\eqref{eq:tree-node-y} for $\eta=\epsilon$. Note that $\mu([0,q])=a(q)$
and $\mu=\nu$ on $(q,1]$. We record the result in the form
\begin{equation}\label{eq:tree-atom-safe}
 \partial_z^2U(q)=1-\int_{[0,1]}\E\bigl[\varpi(q\vee r)^2\bigm|\calG_q\bigr]
   \,\mu(\dd r).
\end{equation}

\paragraph{Step 2: arbitrary profiles.}
Let $\mu_m\Rightarrow\mu$ be the step profiles of the proof of
\cref{lem:tree-regularity}. Applying the proof of \cref{lem:tree-density} to
the single edge from $q$ to $1$ followed by the Gibbs probability, we see
that \eqref{eq:tree-future-density} with
$z=0$ is the density of the future given the field at $q$ for every
profile. Its exponent equals
$\mathcal E(\sigma;Z_1)-\int_{[0,1]}U(q\vee r,Z_{q\vee r})\,\mu(\dd r)$,
and the integrand is continuous in $r$. The reference law does not depend
on the profile, $U_m\to U$ uniformly, and $\mu_m\Rightarrow\mu$. Then
the densities for $\mu_m$ converge pointwise to the density for $\mu$.
Writing $L_m$ for the density for $\mu_m$, we have
\[
 \log L_m=\mathcal E(\sigma;Z_1)-U_m(1,Z_1)
   +\int_{[0,1]}\bigl[U_m(1,Z_1)-U_m(q\vee r,Z_{q\vee r})\bigr]\,
   \mu_m(\dd r),
\]
since $\mu_m([0,1])=1$. The first term is nonpositive, because
$U_m(1,\cdot)$ is the logarithm of a sum that contains
$e^{\mathcal E(\sigma;\cdot)}$. Using this, the bounds on the first derivatives in
\cref{lem:tree-regularity}, and the Lipschitz bound in $q$ from Step~3 of
its proof, we find
\[
 \log L_m\le C_N\Bigl(nc+b+\sup_{q\le r\le1}|Z_r|\Bigr),
\]
where $C_N$ depends only on $N$. The same bound holds for the density for
$\mu$, and the densities have a common Gaussian integrable envelope. They
therefore converge in $L^1$ of the reference law, and expectations of
bounded functions converge. The left side of \eqref{eq:tree-first-derivative} converges by
\cref{lem:tree-regularity}, so $\partial_zU(r)=\E[\eta\mid\calG_r]$ also
for $\mu$. Next, we consider the integrand
$r\mapsto\E[\varpi(q\vee r)^2\mid\calG_q]$ in \eqref{eq:tree-atom-safe}.
Its value for $\mu_m$ differs from its value for $\mu$ by at most the
$L^1$ distance of the two densities plus twice the supremum of the
difference of the first derivatives. It therefore converges uniformly
in $r$.
The limit is continuous in $r$, because the first derivatives are
jointly continuous by \cref{lem:tree-regularity}. Together with the weak
convergence of $\mu_m$, this proves \eqref{eq:tree-atom-safe}, and hence
\eqref{eq:nodeTheta} and \eqref{eq:tree-node-y}, for $\mu$. In these
identities, an atom of $\nu$ at $q$ is counted in $a(q)$, and an atom at
$1$ is counted in the integral.

\paragraph{Step 3: the derivative in $b$ for step profiles.}
Let $\mu$ be a step profile and $q_\ell\le q<q_{\ell+1}$, with $a=a_\ell$
on $[q_\ell,q_{\ell+1})$. Suppose first that $a_\ell>0$, and let
$\phi=e^{a_\ell U}$. On this interval, $\phi$ solves the backward heat
equation $\partial_q\phi+\frac12\Tr(\Lambda_b\nabla^2\phi)=0$. In particular,
$\phi(q)$ is the Gaussian average of $\phi(q_{\ell+1})$ over an increment
of covariance $(q_{\ell+1}-q)\Lambda_b$. We differentiate in $b$, using
the Gaussian heat identity in the last coordinate, and divide by
$a_\ell\phi(q)$. We obtain
\[
 \partial_bU(q)=\E\bigl[\partial_bU(q_{\ell+1})\bigm|\calG_q\bigr]
  +\frac{q_{\ell+1}-q}2\,
  \E\bigl[(\partial_{yy}U+a_\ell(\partial_yU)^2)(q_{\ell+1})\bigm|\calG_q\bigr].
\]
The function
$\partial_{yy}\phi=a_\ell(\partial_{yy}U+a_\ell(\partial_yU)^2)\phi$ also
solves the backward heat equation. Further, the transition of the field
from $q$ to $r\le q_{\ell+1}$ is the Gaussian transition tilted by
$\phi(r)/\phi(q)$. It follows that
$\E[(\partial_{yy}U+a_\ell(\partial_yU)^2)(r)\mid\calG_q]$ does not
depend on $r\in[q,q_{\ell+1}]$, and that the last term equals
$\frac12\int_q^{q_{\ell+1}}\E[(\partial_{yy}U+a(\partial_yU)^2)(r)\mid\calG_q]\dd r$.
The case $a_\ell=0$ is the same with $\phi=U$. Iterating over the grid
intervals and using $\partial_bU(1)=0$, we find
\begin{equation}\label{eq:tree-b-response}
 \partial_bU(q)=\frac12\int_q^1
   \E\bigl[\partial_{yy}U(r)+a(r)\,m_N(r)^2\bigm|\calG_q\bigr]\dd r,
\end{equation}
where we used that $\partial_yU(r)=m_N(r)$.

\paragraph{Step 4: the derivative in $b$ for arbitrary profiles.}
By \eqref{eq:tree-node-y}, the integrand in \eqref{eq:tree-b-response}
lies in $[0,1]$. Let $R(b)$ be the right side of
\eqref{eq:tree-b-response} for $\mu$. For the step profiles $\mu_m$, the
right side of \eqref{eq:tree-b-response} converges to $R(b)$ at each
fixed $b$. Indeed, the densities converge as in Step 2, the derivatives
converge uniformly, and $a_m\to a$ in $L^1(0,1)$. These expressions are
bounded by $1/2$. By dominated convergence in the identity
$U_m(b_2)-U_m(b_1)=\int_{b_1}^{b_2}\partial_bU_m\dd b$, we obtain
$U(b_2)-U(b_1)=\int_{b_1}^{b_2}R(b)\dd b$. We write the last-field
increments of the reference law as $\sqrt b$ times a fixed Brownian
motion. Using \cref{lem:tree-regularity} and dominated convergence, we then see
that, for each fixed $q$, $R$ is jointly continuous in
$(b,\theta,x,y)$. It follows that
$\partial_bU$ exists, equals $R$, and satisfies
\eqref{eq:tree-b-response} for $\mu$. Inserting \eqref{eq:tree-node-y}
at level $r$ and using the tower property, we have
\[
 \partial_bU(q)=\frac12\int_q^1\Bigl(1-\int_{(r,1]}
   \E[m_N(\ell)^2\mid\calG_q]\,\nu(\dd\ell)\Bigr)\dd r .
\]
By Fubini's theorem, we have $\int_q^1\dd r\int_{(r,1]}F(\ell)\,\nu(\dd\ell)
=\int_{(q,1]}(\ell-q)F(\ell)\,\nu(\dd\ell)$ for bounded $F\ge0$, and
\eqref{eq:nodeb} follows. Finally, the integral in \eqref{eq:nodeb} lies
between $0$ and $(1-q)\nu((q,1])\le1-q$, which gives
$0\le\partial_bU(q)\le(1-q)/2$.
\end{proof}

\subsection{Differentiation in the interpolation parameter}

For a point $v$ of $\mathcal T$, we write $\hat v$ and $\check v$ for new
leaves whose paths leave the path through $v$ at the level $q_v$. If $v$
is an interior point of an edge, the extended tree has a new vertex at
$v$. If $v$ is a vertex, the new edge starts at $v$, and if $v$ is a
terminal point, the new leaf is attached there. Distinct new leaves are
conditionally independent, also when they are attached at the same
point. For $r>q_v$, we write $\hat r$ and $\check r$ for two new leaves
that follow a common new branch from $v$ to the level $r$ and split
there. Given the couplings and the field of $\mathcal T$, a new branch
from $v$ is a continuation of the diffusion from the field at $v$, and it
is independent of everything else. Let $F$ be a bounded function of the
spins of $\mathcal T$, and let $G,G'$ be bounded functions of one spin.
Then
\begin{equation}\label{eq:tree-replacement}
 \begin{aligned}
 \E\bigl[F\,G(\sigma^{\hat v})\bigr]
 &=\E_{\mathcal T,u}\bigl[F\,\E[G(\sigma)\mid\calG_{q_v}](v)\bigr],\\
 \E\bigl[F\,G(\sigma^{\hat r})G'(\sigma^{\check r})\bigr]
 &=\E_{\mathcal T,u}\Bigl[F\,\E\bigl[\E[G(\sigma)\mid\calG_r]\,
   \E[G'(\sigma)\mid\calG_r]\bigm|\calG_{q_v}\bigr](v)\Bigr].
 \end{aligned}
\end{equation}
Here the left sides are taken under the extended trees. A conditional
expectation under $\Pp_u$ given $\calG_{q_v}$ is a function of the
couplings and of the field at level $q_v$, and $(v)$ means that it is
evaluated at the couplings and at the field at $v$. In the same way, two
new leaves $\hat v$ and $\check w$ attached at $v$ and $w$ represent the
product of the two conditional expectations, evaluated at $v$ and at $w$.

Let $\varphi_{\mathsf{ab}}$ be a quantity attached to each pair of
distinct leaves of an extended tree. For a bounded function $F$ of the
spins of $\mathcal T$, we define the five-pattern sum
\[
 \begin{aligned}
 \mathfrak S_\varphi(F)={}&
 \sum_{\mathsf a\ne\mathsf b}\E(F\epsilon_{\mathsf a}\epsilon_{\mathsf b}
   \varphi_{\mathsf{ab}})
 -2\sum_{\mathsf a,v}\omega_v\E(F\epsilon_{\mathsf a}\epsilon_{\hat v}
   \varphi_{\mathsf a\hat v})
 +\sum_{v,w}\omega_v\omega_w\E(F\epsilon_{\hat v}\epsilon_{\check w}
   \varphi_{\hat v\check w})\\
 &+\sum_v\omega_va(q_v)\E(F\epsilon_{\hat v}\epsilon_{\check v}
   \varphi_{\hat v\check v})
 +\sum_v\omega_v\int_{(q_v,1]}\E(F\epsilon_{\hat r}\epsilon_{\check r}
   \varphi_{\hat r\check r})\,\nu(\dd r).
 \end{aligned}
\]
The sums run over leaves $\mathsf a,\mathsf b$ of $\mathcal T$ and points
$v,w$ of $\mathcal T$, and the weights $\omega$ are those of
$\mathcal T$. We take each expectation under the extended tree that
contains the leaves involved. We use this sum with
$\varphi_{\mathsf{ab}}=R^-_{\mathsf{ab}}$, with
$\varphi_{\mathsf{ab}}=q_{\mathsf{ab}}$, and with
$\varphi_{\mathsf{ab}}=\Delta_{\mathsf{ab}}$. Then
$\mathfrak S_\Delta=\mathfrak S_{R^-}-\mathfrak S_q$, and
$\epsilon_{\mathsf a}\epsilon_{\mathsf b}\Delta_{\mathsf{ab}}=D_{\mathsf{ab}}$.
The following table lists the five patterns. Its last column gives the
level $q_{\mathsf{ab}}$ of the pair.
\begin{center}
\small
\begin{tabular}{@{}llll@{}}
\toprule
pair of leaves & new leaves & coefficient & level of the pair\\
\midrule
two leaves $\mathsf a\ne\mathsf b$ of $\mathcal T$ & none & $1$
  & $q_{\mathsf{ab}}$\\
a leaf $\mathsf a$ and a leaf attached at $v$ & $\hat v$
  & $-2\,\omega_v$ & $q_{\mathsf a\hat v}$\\
leaves attached at $v$ and at $w$ & $\hat v,\check w$
  & $\omega_v\omega_w$ & $q_{\hat v\check w}$\\
two leaves attached at $v$ & $\hat v,\check v$ & $a(q_v)\,\omega_v$
  & $q_v$\\
two leaves splitting at $r>q_v$ & $\hat r,\check r$
  & $\omega_v\,\nu(\dd r)$ & $r$\\
\bottomrule
\end{tabular}
\end{center}

\begin{lemma}[Differentiation under the tree law]\label{lem:streaming}
Let $\mu$ be a deterministic profile, $\mathcal T$ a marked tree with $k$
leaves, and $F$ a bounded function of the spins of its leaves. Then
$u\mapsto\E_{\mathcal T,u}F$ belongs to $C^2([0,1])$, with one-sided
derivatives at the endpoints, and
\begin{align}
 \frac{\dd}{\dd u}\E_{\mathcal T,u}F=\frac\alpha2\bigg[&
 \sum_{\mathsf a\ne\mathsf b}\E(FD_{\mathsf{ab}})
 -2\sum_{\mathsf a,v}\omega_v\E(FD_{\mathsf a\hat v})
 +\sum_{v,w}\omega_v\omega_w\E(FD_{\hat v\check w})\nonumber\\
 &+\sum_v\omega_va(q_v)\E(FD_{\hat v\check v})
 +\sum_v\omega_v\int_{(q_v,1]}\E(FD_{\hat r\check r})\,\nu(\dd r)
 \nonumber\\
 &+s\,\E_{\mathcal T,u}F\int\E_{2,r,u}D_{12}\,\nu(\dd r)\bigg].
 \label{eq:streaming}
\end{align}
The sums run over leaves $\mathsf a,\mathsf b$ and points $v,w$ of
$\mathcal T$, with the weighted sums interpreted as integrals against
$\omega$ as above. The first five expectations are taken under extended
trees, and the last term is a product of expectations over independent
copies. The absolute coefficients in the bracket sum to at most
$4k^2+2k+1$. The second derivative is obtained by applying the same
formula to each tree moment on its right side, and the product rule to
independent products of moments.
\end{lemma}

In the notation above, \eqref{eq:streaming} reads
$\frac{\dd}{\dd u}\E_{\mathcal T,u}F=\frac\alpha2[\mathfrak S_\Delta(F)
+s\,\E_{\mathcal T,u}F\int\E_{2,r,u}D_{12}\,\nu(\dd r)]$.

\begin{proof}
If $\alpha=0$, the law does not depend on $u$, and both sides vanish. Let
$\alpha>0$. The proof has five steps. We check
that \cref{lem:fnd-path-heat} applies, and we compute the contributions
of the row and of the last field to the derivative of the unnormalized
moment. We then compute the derivative of the normalizer. Finally, we
bound the coefficients and prove the $C^2$ regularity.

\paragraph{Step 1: the heat identity.}
By \cref{lem:tree-density}, we have
$\E_{\mathcal T,u}F=Z_s(u)^{-1}\E_{\rm ref}[F\mathscr D_{\mathcal T}]$,
where $\E_{\rm ref}$ denotes integration against the reference measure of
that lemma. We write the row as $\theta=(\alpha u/n)^{1/2}(g_{iN})_{i\le n}$,
and we use \cref{lem:fnd-path-heat} on $I=[0,1]$ with the following
data. The finite coordinates are the entries of $\theta$, with variances
$\alpha u/n$. The field paths are the $n+1$ coordinates of the driftless
field on the edges of $\mathcal T$. For an edge $e$ with endpoint levels
$r_e<h_e$, we use the local coordinate $v\in[0,1]$, and the clocks are
\[
 T_{e,i,u}(v)=c\,(h_e-r_e)\,v\quad(i\le n),\qquad
 T_{e,N,u}(v)=b(u)\,(h_e-r_e)\,v .
\]
The couplings $g_{ij}$, $i<j\le n$, do not depend on $u$, and the spins
are the counting variable. The row, the Brownian motions of the
different edges and components, and the couplings $g_{ij}$ are mutually
independent. We take $\ell=\log\mathscr D_{\mathcal T}$ and
$\Phi=F$. The elements of the space $\mathcal C$ of
\cref{lem:fnd-path-heat} need not agree at the common
endpoints of edges. To define $\ell$ on the whole space, we evaluate the field at each nonroot vertex at the
endpoint of its incoming edge. At the root, we choose a fixed outgoing
edge $e_0$ and evaluate the field at the initial point of $e_0$; its
coordinates vanish on every path of the driftless field. This agrees with
\eqref{eq:treedensity} on the paths of the driftless field. The
derivative measures below, including the atom of $\omega$ at the root,
are those of this extension. We check the hypotheses \textup{(A1)}--\textup{(A4)} of
\cref{lem:fnd-path-heat}.
\begin{itemize}
\item \textup{(A1)}: the variances $\alpha u/n$ are positive for $u>0$.
The clocks are affine in $u$, and they are strictly increasing in $v$,
because $c>0$, $b(u)\ge c$, and every edge has positive length.
\item \textup{(A2)}: the covariance of the last component at two field
points $v,w$ is $b(u)h(v,w)$, where $h(v,w)$ is the level of the last
common point of the paths from the root to $v$ and to $w$. For the other
components, it is $c\,h(v,w)$. Hence $K$ is affine in $u$, with
$\dot K=(\alpha/n)I_n$ on the row block, $\dot K=-\alpha h$ on the
last-field block, and $\dot K=0$ elsewhere. These kernels are
continuous.
\item \textup{(A3)}: $\Phi=F$ is bounded and $\dot\Phi=0$. The explicit
derivative of $\ell$ comes from $b(u)$ in $U$, so that
$\dot\ell=-b'(u)\sum_v\omega_v\partial_bU(v)=\alpha\sum_v\omega_v
\partial_bU(v)$. The levels $q_v$ and the finite measure $\omega$ are
fixed. Using the fixed-level continuity and uniform bound of
\cref{lem:tree-responses}, together with dominated convergence against $\omega$, we see
that $\dot\ell$ is continuous in $(u,G)$; it is bounded by $\alpha k/2$.
If we treat each rest coupling $g_{ij}$ as an additional source parameter
with bounded feature $\beta\sqrt{t/N}\,\tau_i\tau_j$ in Step~1 of the proof of
\cref{lem:tree-regularity}, we obtain joint continuity of $U$ and its spatial
derivatives also in the rest couplings. The fixed-level response formula
of \cref{lem:tree-responses} then gives joint continuity of $\partial_bU$
in those couplings by dominated convergence, using the uniform bound of
that lemma.
Thus $\ell$ and $\dot\ell$ are jointly continuous in $(u,G,g)$, as required
by \textup{(A3)}.
The first derivatives of $U$ in
the row and in the field are bounded by one. Its derivative in $g_{ij}$,
$i<j\le n$, is the conditional mean of $\beta\sqrt{t/N}\,\tau_i\tau_j$,
by the argument of Step~1 of the proof of \cref{lem:tree-regularity}.
Integrating \eqref{eq:cavityPDE} in $q$ and using that the first spatial
derivatives of $U$ are bounded by one and the second by two (Step~1 of that
proof), we have $|U(q,0,0)|\le N\log2+\frac32(nc+b(u))$
when all couplings, the row, and the field vanish.
Since $\mathcal E$ is linear in the couplings, the row, and the field, we
conclude that $|\ell|\le C(1+|g|+|\theta|+\|G^{\rm path}\|_\infty)$, where
$C$ depends only on $N$, $k$, and $\beta$. Here $g$ collects the couplings
$g_{ij}$, $i<j\le n$.
\item \textup{(A4)}: the score of $\ell$ consists of the atoms
$\tau_i^{\mathsf a}\epsilon_{\mathsf a}$ (row), $\tau^{\mathsf a}$ and
$\epsilon_{\mathsf a}$ (field at the terminal point of $\mathsf a$), and
of $-\nabla U(v)\,\omega(\dd v)$, where $\nabla$ is the gradient in the
row and the field. Since $\mathcal E$ is linear in these variables, the
Hessian of $\ell$ is $-\nabla^2U(v)\,\omega(\dd v)$, placed at the pairs
of indices of the point $v$. It has a row--row block
$-\partial_{\theta_i\theta_j}U(v)\,\omega(\dd v)$, mixed row--field
blocks such as $-\partial_{\theta_iy}U(v)\,\omega(\dd v)$ and
$-\partial_{\theta_ix_j}U(v)\,\omega(\dd v)$ with their transposes, and
field--field blocks such as $-\partial_{yy}U(v)\,\omega(\dd v)$, placed at
the diagonal pairs of index points of $v$. By
\cref{lem:tree-regularity}, these derivatives of $U$ are bounded and
jointly continuous in $(b,q,\theta,x,y)$. As the points $(q_v,X_v,Y_v)$
stay in a compact set when $G$ varies in a bounded set, the measures are
continuous in total variation in $(u,G)$. Their total variations are
at most $C_Nk$. The measures for $\Phi$ vanish.
\end{itemize}
The constants in these hypotheses depend on $N$ and $k$. We use them
only to justify the differentiation, and the identity below has
coefficients that do not depend on $N$. Since $\Phi=F$ depends neither on
the Gaussian coordinates nor on $u$, \eqref{eq:fnd-path-heat} and
\eqref{eq:fnd-path-normalizer} imply that
\begin{equation}\label{eq:tree-quotient}
 \frac{\dd}{\dd u}\E_{\mathcal T,u}F
 =\frac{\partial_u\E_{\rm ref}[F\mathscr D_{\mathcal T}]}{Z_s(u)}
 -\E_{\mathcal T,u}F\,\frac{\dd}{\dd u}\log Z_s(u),
\end{equation}
where the derivative of the unnormalized moment is
\[
 \frac{\partial_u\E_{\rm ref}[F\mathscr D_{\mathcal T}]}{Z_s(u)}
 =\E_{\mathcal T,u}\Bigl[F\Bigl(\dot\ell+\frac\alpha{2n}\sum_{i=1}^n
   \bigl((\partial_{\theta_i}\ell)^2+\partial_{\theta_i}^2\ell\bigr)
   -\frac\alpha2\bigl\langle h,S_y\otimes S_y+M_{yy}\bigr\rangle\Bigr)\Bigr].
\]
Here $S_y$ and $M_{yy}$ are the score and Hessian measures of $\ell$ in
the last-field coordinates.

\paragraph{Step 2: the row.}
By \eqref{eq:treedensity} and \eqref{eq:nodeTheta}, we have
\[
 \partial_{\theta_i}\ell=\sum_{\mathsf a}\tau_i^{\mathsf a}\epsilon_{\mathsf a}
   -\sum_v\omega_v\pi_i(v),
 \qquad
 \partial_{\theta_i}^2\ell=-\sum_v\omega_v\partial_{\theta_i}^2U(v),
\]
where $\pi_i(v)$ and $\partial_{\theta_i}^2U(v)$ are evaluated at the
couplings and at the field at $v$. We expand the square and insert
\eqref{eq:nodeTheta}. The diagonal terms
$(\tau_i^{\mathsf a}\epsilon_{\mathsf a})^2=1$ contribute $k$, and the
constant $1$ in \eqref{eq:nodeTheta} contributes $-\omega(\mathcal T)$.
By \cref{lem:tree-density}, these constants add up to $s$. The other
terms of $(\partial_{\theta_i}\ell)^2+\partial_{\theta_i}^2\ell$ are
\[
 \begin{aligned}
 &\sum_{\mathsf a\ne\mathsf b}\tau_i^{\mathsf a}\epsilon_{\mathsf a}
   \tau_i^{\mathsf b}\epsilon_{\mathsf b}
 -2\sum_{\mathsf a,v}\omega_v\tau_i^{\mathsf a}\epsilon_{\mathsf a}\pi_i(v)
 +\sum_{v,w}\omega_v\omega_w\pi_i(v)\pi_i(w)\\
 &\quad+\sum_v\omega_va(q_v)\pi_i(v)^2
 +\sum_v\omega_v\int_{(q_v,1]}\E[\pi_i(r)^2\mid\calG_{q_v}](v)\,\nu(\dd r).
 \end{aligned}
\]
By \eqref{eq:tree-replacement}, inside $\E_{\mathcal T,u}[F\,\cdot\,]$
we may replace $\pi_i(v)$ by $\tau_i^{\hat v}\epsilon_{\hat v}$, and
$\pi_i(v)\pi_i(w)$ by the product for $\hat v$ and $\check w$. We may
also replace $\pi_i(v)^2$ by the product for $\hat v$ and $\check v$, and
$\E[\pi_i(r)^2\mid\calG_{q_v}](v)$ by the product for $\hat r$ and
$\check r$. Averaging over $i$ then replaces
$\tau_i^{\mathsf a}\tau_i^{\mathsf b}$ by $R^-_{\mathsf{ab}}$. We obtain
\begin{equation}\label{eq:tree-row-part}
 \frac\alpha{2n}\sum_{i=1}^n\E_{\mathcal T,u}\bigl[F\bigl(
   (\partial_{\theta_i}\ell)^2+\partial_{\theta_i}^2\ell\bigr)\bigr]
 =\frac\alpha2\bigl[s\,\E_{\mathcal T,u}F+\mathfrak S_{R^-}(F)\bigr].
\end{equation}

\paragraph{Step 3: the last field.}
By \eqref{eq:treedensity} and \eqref{eq:tree-node-y}, the score and
Hessian measures of $\ell$ in the last-field coordinates are
\[
 S_y=\sum_{\mathsf a}\epsilon_{\mathsf a}\delta_{\mathsf a}
                      -m_N(v)\,\omega(\dd v),\qquad
 M_{yy}=-\operatorname{diag}_{\#}
                  \bigl(\partial_{yy}U(v)\,\omega(\dd v)\bigr),
\]
where $\delta_{\mathsf a}$ is the unit mass at the terminal point of
$\mathsf a$. Here $\operatorname{diag}_{\#}$ is the image of a measure
under $v\mapsto(v,v)$, so that atoms of $\omega$ stay at exact diagonal
pairs. Let a leaf stand for its terminal point. Then the kernel $h$
satisfies $h(\mathsf a,\mathsf a)=1$,
$h(\mathsf a,\mathsf b)=q_{\mathsf{ab}}$,
$h(\mathsf a,v)=q_{\mathsf a\hat v}$, $h(v,w)=q_{\hat v\check w}$, and
$h(v,v)=q_v$. Expanding $\langle h,S_y\otimes S_y\rangle$ as in Step 2
and replacing $m_N(v)$ by $\epsilon_{\hat v}$ with
\eqref{eq:tree-replacement}, we find
\[
 \begin{aligned}
 -\frac\alpha2\E_{\mathcal T,u}\bigl[F\langle h,S_y\otimes S_y\rangle\bigr]
 =-\frac\alpha2\Bigl[&k\,\E_{\mathcal T,u}F
 +\sum_{\mathsf a\ne\mathsf b}q_{\mathsf{ab}}
   \E(F\epsilon_{\mathsf a}\epsilon_{\mathsf b})\\
 &-2\sum_{\mathsf a,v}\omega_vq_{\mathsf a\hat v}
   \E(F\epsilon_{\mathsf a}\epsilon_{\hat v})
 +\sum_{v,w}\omega_v\omega_wq_{\hat v\check w}
   \E(F\epsilon_{\hat v}\epsilon_{\check w})\Bigr].
 \end{aligned}
\]
The bracket consists of the constant $k\,\E_{\mathcal T,u}F$ and the
first three patterns of $\mathfrak S_q(F)$. The Hessian measure
contributes
$-\frac\alpha2\langle h,M_{yy}\rangle=\frac\alpha2\sum_v\omega_vq_v
\partial_{yy}U(v)$, and we recall that
$\dot\ell=\alpha\sum_v\omega_v\partial_bU(v)$. Together, these two terms
give $\frac\alpha2\sum_v\omega_v[q_v\partial_{yy}U(v)+2\partial_bU(v)]$.
By \eqref{eq:tree-node-y} and \eqref{eq:nodeb}, we have
\[
 q_v\partial_{yy}U(v)+2\partial_bU(v)
 =1-q_va(q_v)m_N(v)^2
   -\int_{(q_v,1]}r\,\E[m_N(r)^2\mid\calG_{q_v}](v)\,\nu(\dd r).
\]
In the integral, we combined $q_v\int_{(q_v,1]}$ from
\eqref{eq:tree-node-y} with $\int_{(q_v,1]}(r-q_v)$ from \eqref{eq:nodeb}.
After the replacements \eqref{eq:tree-replacement}, the constant gives
$\frac\alpha2\omega(\mathcal T)\E_{\mathcal T,u}F
=\frac\alpha2(k-s)\E_{\mathcal T,u}F$. The other two terms are the last
two patterns of $-\frac\alpha2\mathfrak S_q(F)$, with the levels
$q_{\hat v\check v}=q_v$ and $q_{\hat r\check r}=r$. Altogether, the last
field and the explicit derivative contribute
$\frac\alpha2[-k+(k-s)]\E_{\mathcal T,u}F-\frac\alpha2\mathfrak S_q(F)$.
Adding \eqref{eq:tree-row-part}, we see that the constants $\alpha s/2$,
$-\alpha k/2$, and $\alpha(k-s)/2$ cancel, and since
$\mathfrak S_{R^-}-\mathfrak S_q=\mathfrak S_\Delta$, we conclude that
\begin{equation}\label{eq:tree-unnormalized}
 \frac{\partial_u\E_{\rm ref}[F\mathscr D_{\mathcal T}]}{Z_s(u)}
 =\frac\alpha2\,\mathfrak S_\Delta(F).
\end{equation}
In words, the row inserts a rest overlap, the last field inserts the
level of the corresponding pair, and their difference is an overlap
defect.

\paragraph{Step 4: the normalizer.}
We compute the derivative of $Z_s(u)=\E_{\rm ref}e^{sU(0,0,0)}$ with
\eqref{eq:fnd-path-normalizer}; now the only active coordinates are those of the
row. The row contributes
$\frac{\alpha s}{2n}\sum_i\E_u[\partial_{\theta_i}^2U(0)
+s(\partial_{\theta_i}U(0))^2]$. As $a(0)-s=\nu(\{0\})$, we obtain from
\eqref{eq:nodeTheta} that
\[
 \partial_{\theta_i}^2U(0)+s\bigl(\partial_{\theta_i}U(0)\bigr)^2
 =1-\int_{[0,1]}\E[\pi_i(r)^2\mid\calG_0]\,\nu(\dd r).
\]
Here, as in \cite[(2.11)]{AL}, the term $s(\partial_{\theta_i}U)^2$
removes the prescribed root mass $s$, and a possible atom $\nu(\{0\})$
remains in the integral. Since
$\E_u\E[\pi_i(r)^2\mid\calG_0]=\E_{2,r,u}[\tau_i^1\epsilon_1\tau_i^2\epsilon_2]$,
the row contributes
$\frac{\alpha s}2[1-\int\E_{2,r,u}[\epsilon_1\epsilon_2R^-_{12}]\,
\nu(\dd r)]$. The coefficient $b(u)$ contributes
$sb'(u)\E_u\partial_bU(0)=-\alpha s\E_u\partial_bU(0)$. By
\eqref{eq:nodeb}, this equals
$-\frac{\alpha s}2[1-\int r\,\E_{2,r,u}[\epsilon_1\epsilon_2]\,\nu(\dd r)]$.
Adding the two contributions, we find
\begin{equation}\label{eq:normalizerderivative}
 \frac\dd{\dd u}\log Z_s(u)
 =-\frac{\alpha s}{2}\int\E_{2,r,u}D_{12}\,\nu(\dd r).
\end{equation}
Inserting \eqref{eq:tree-unnormalized} and
\eqref{eq:normalizerderivative} into \eqref{eq:tree-quotient}, we obtain
\eqref{eq:streaming}. In its last term, $\E_{2,r,u}$ is the expectation
under an independent copy of the complete tree law, including its
couplings.

\paragraph{Step 5: coefficients and regularity.}
Recall that $\omega(\mathcal T)=k-s\le k$, $a\le1$, and $\nu([0,1])\le1$. The
absolute coefficients of the six terms in \eqref{eq:streaming} are
at most $k(k-1)$, $2k^2$, $k^2$, $k$, $k$, and $1$. Their sum is
$4k^2+k+1\le4k^2+2k+1$.

We now show that each extended-tree moment on the right side of
\eqref{eq:streaming} is continuous in $u$. We write the row as
$(\alpha u/n)^{1/2}g$, and the last-field coordinates as $\sqrt{b(u)}$
times a fixed Brownian field on the tree. By \cref{lem:tree-regularity},
the density is then continuous in $u$ for fixed Gaussian variables. Its
linear growth gives a Gaussian envelope that is uniform in $u$. By
dominated convergence, the moments and their normalizers are continuous.
The weights and the measures $\omega$ and $\nu$ over which the right side
is integrated do not depend on $u$, and the integrands are bounded by
$2\|F\|_\infty$. Bounded convergence then shows that the right side is
continuous. The derivative exists at every $u\in[0,1]$, one-sided at the
endpoints, so $u\mapsto\E_{\mathcal T,u}F$ is $C^1$.
Next, we apply \eqref{eq:streaming} to each bounded spin observable $FD$
on the extended trees; these observables do not depend on $u$. The
numbers of leaves grow by at most two, and the coefficients are bounded
in terms of $k$. Each of these moments therefore has a derivative bounded
by $C_k\alpha\|F\|_\infty$, uniformly in the points $v,w$ and the level
$r$. This lets us differentiate under the integrals against $\omega$ and
$\nu$ by dominated convergence, and we use the product rule for the last
term. The resulting right side is continuous by the same argument. This
proves the $C^2$ regularity and the rule for the second derivative.
\end{proof}

\begin{corollary}[Comparison of nonnegative observables]\label{cor:comparison}
For every deterministic profile, every marked tree $\mathcal T$ with $k$
leaves, and every nonnegative bounded function $F$ of the spins of its
leaves,
\begin{equation}\label{eq:comparison}
 \E_{\mathcal T,u}F\le K_{\beta,k}\E_{\mathcal T,v}F
 \quad(0\le u,v\le1),
 \qquad K_{\beta,k}=\exp\bigl(\beta^2(4k^2+2k+1)\bigr).
\end{equation}
\end{corollary}

\begin{proof}
Since $|D_{\mathsf{ab}}|\le2$ and $F\ge0$, each expectation in
\eqref{eq:streaming} is at most twice the expectation of $F$ under the
corresponding extended tree. By deletion, the latter equals
$\E_{\mathcal T,u}F$, and the same bound holds for the product in the
last term. Using \cref{lem:streaming} and $\alpha\le\beta^2$, we find
\[
 \bigl|\partial_u\E_{\mathcal T,u}F\bigr|
 \le\beta^2(4k^2+2k+1)\,\E_{\mathcal T,u}F,\qquad0\le u\le1 .
\]
By Gronwall's inequality on the interval between $u$ and $v$, we obtain
$\E_{\mathcal T,u}F\le e^{\beta^2(4k^2+2k+1)|u-v|}\E_{\mathcal T,v}F$.
In particular, \eqref{eq:comparison} holds, also in the case
$\E_{\mathcal T,v}F=0$, in which the moment vanishes for all $u$. The
comparisons in \cite[Lemmas~2.9 and~4.3]{OpenAIFluctuationScale} are
proved by the same differential inequality. The classical analog
on one site is \cite[Lemma~1.6.4]{TalagrandVolI2011}.
\end{proof}

\section{Products of three overlap defects}\label{subsec:c7-three}

In the Taylor remainder of the cavity covariance, and in the variance
comparison of \cref{lem:cl-star}, we differentiate tree moments of
observables that already contain overlap defects. Each differentiation
adds one more defect. We show that integrated moments of products of
three defects are controlled by the integrated third absolute moment
\begin{equation}\label{eq:cav-M3}
 \mathcal M_3(u)=\int\E_{2,q,u}|R^-_{12}-q|^3\,\nu(\dd q)
\end{equation}
of the overlap defect of two descendants. In \cref{lem:c7-thermal},
we bound $\mathcal M_3$ by $A$, $J$, and $H_3$.

\paragraph{Terms and marks.}
Fix integers $k_0\ge1$ and $j\ge0$. For each $p\in[0,1]$, let
$\mathcal T_p$ be the tree of $k_0$ leaves that split at $p$, and let
$F_p$ be a bounded function of their spins, jointly measurable in $p$ and
the spins. We consider the integrated tree moment
$\int\nu(\dd p)\E_{\mathcal T_p,u}F_p$. We also consider the terms
produced from it by $j$ applications of \eqref{eq:streaming}, with the
product rule for products of moments over independent trees. Each
application multiplies by $\alpha/2$ and replaces a moment by the six
terms of \eqref{eq:streaming}. Each term produced after $j$ applications
is therefore $(\alpha/2)^j$ times an integral against a signed measure on
finitely many \emph{marks}, which are the following variables. The marks
are $p$, the points $v$ and $w$ at which the differentiations attached
leaves, the levels $r$ at which new branches split, and the split levels
$r$ of the independent trees produced by the last term of
\eqref{eq:streaming}. The integrand is a product of expectations over
independent trees, whose shapes are determined by the marks. The signed
measure collects the coefficients $1$, $-2\omega_v$, $\omega_v\omega_w$,
$a(q_v)\omega_v$, $\omega_v\nu(\dd r)$, and $s\,\nu(\dd r)$ of the
attachments, and the multiplicities of the product rule. We let
$\mathfrak m$ be its total variation. The factor $(\alpha/2)^j$ is not
part of $\mathfrak m$. The level of a mark is $p$, $q_v$, $q_w$, or $r$.

\begin{lemma}[Three-defect bound]\label{lem:cav-threeproducts}
For every deterministic profile, every such term containing three
defects $\Delta_1,\Delta_2,\Delta_3$, of pairs of leaves in the same or
in independent trees, satisfies
\begin{equation}\label{eq:threeproducts}
 \int\E\bigl|\Delta_1\Delta_2\Delta_3\bigr|\,\dd\mathfrak m
 \le C\mathcal M_3(u).
\end{equation}
Here $\E$ is the expectation under the product of the tree laws of the
term, and $C$ depends only on the numbers of leaves and of
differentiations.
\end{lemma}

The bounds on the marks in Step 1 of the proof are a counterpart, for an
arbitrary profile, of the bounds on split levels in
\cite[Lemma~2.8]{OpenAIFluctuationScale}.

\begin{proof}
The proof has three steps. We show that the level of every mark has a
marginal under $\mathfrak m$ dominated by a multiple of $\nu$, and that
the level of every pair of leaves is the level of a mark. We then use
Hölder's inequality.

\paragraph{Step 1: marginals of the marks.}
Each differentiation adds at most two leaves to one tree, or creates an
independent tree with two leaves. After $i$ differentiations, every tree
of a term has at most $K_i=k_0+2i$ leaves, and a product has at
most $k_0+i$ factors. Let $\mathfrak m_i$ be the measure $\mathfrak m$ of
a term after $i\le j$ differentiations. We show by induction on $i$ that
$\mathfrak m_i$ has total mass at most $C_i$, and that the level of each
of its marks has a marginal at most $C_i\nu$. Here $C_i$ depends only on
$k_0$ and $i$. For $i=0$, the only mark is $p$, and
$\mathfrak m_0=\nu(\dd p)$.

For the induction step, let $K=K_i$. We first bound the image of the
weights $\omega$ of a tree with at most $K$ leaves under the level map
$v\mapsto q_v$. At most $K$ points of the tree lie at any given level.
The copies of $\nu$ on the edges and the parts $\nu(\{q_v\})$ of the
vertex weights contribute at most $K\nu$. The parts
$(l_v-1)a(q_v)\le l_v-1$ contribute at most
$\sum_v(l_v-1)\delta_{q_v}$, where $v$ runs over the vertices with
$l_v\ge2$, and $\sum_v(l_v-1)\le K$. Note that every such $q_v$ is the
level of an earlier mark. We now consider one more differentiation. By
\cref{lem:streaming}, it multiplies $\mathfrak m_i$ by a kernel of total
mass at most $(k_0+i)(4K^2+2K+1)$; here $k_0+i$ bounds the
number of factors to which the product rule applies. The total mass and
the marginals of the old marks thus grow by at most the factor
$(k_0+i)(4K^2+2K+1)$. A new
point $v$ or $w$ chosen according to $\omega$ has a level whose
conditional law is at most $K[\nu+\sum_{r'}\delta_{r'}]$, where $r'$ runs
over the levels of the old marks. A new level $r$ chosen according to
$\nu(\dd r)$ on $(q_v,1]$, or according to $s\,\nu(\dd r)$, has
conditional law at most $\nu$. Factors $a(q_v)\le1$ and the restriction
$r>q_v$ only decrease these measures. For every Borel set $B$, the
marginal of the level $r_{\rm new}$ of a new mark therefore satisfies
\[
 \mathfrak m_{i+1}(r_{\rm new}\in B)
 \le C\Bigl[\mathfrak m_i(1)\,\nu(B)+\sum_{r'}\mathfrak m_i(r'\in B)\Bigr]
 \le C_{i+1}\,\nu(B),
\]
where $C$ depends only on $k_0$ and $i$. This completes the induction.

\paragraph{Step 2: levels of pairs.}
In a tree of such a term, every vertex with at least two outgoing edges or
leaves lies at the level of a mark. This level is the original split
level $p$, a level $q_v$ at which leaves were attached, a level $r$ at
which a new branch split, or the split level of an independent tree. For
two leaves $\mathsf a,\mathsf b$ of one tree, $q_{\mathsf{ab}}$ is the
level of such a vertex. In particular,
$\chi(q_{\mathsf{ab}})\le\sum_{r'}\chi(r')$ for every $\chi\ge0$, where
$r'$ runs over the levels of the marks. By Step 1, this implies that
\[
 \int\chi(q_{\mathsf{ab}})\,\dd\mathfrak m\le C\int\chi\,\dd\nu .
\]

\paragraph{Step 3: Hölder's inequality.}
As in \cite[Remark~2.4]{OpenAIFluctuationScale}, we apply Hölder's
inequality unconditionally under the positive tree laws, and we bound each
factor after deleting all leaves except one pair. On the product of the
probability space of the term and the mark space with $\mathfrak m$, we
have
\[
 \int\E|\Delta_1\Delta_2\Delta_3|\,\dd\mathfrak m
 \le\prod_{i=1}^3\Bigl(\int\E|\Delta_i|^3\,\dd\mathfrak m\Bigr)^{1/3}.
\]
If we delete all leaves except the two leaves of $\Delta_i$, two
descendants remain, splitting at their common level $q_i$. By deletion, we
have $\E|\Delta_i|^3=\E_{2,q_i,u}|R^-_{12}-q_i|^3$; for a product over
independent trees, the other trees integrate to one. We use Step 2 with
$\chi(q)=\E_{2,q,u}|R^-_{12}-q|^3$, and we find that each factor is at
most $(C\mathcal M_3(u))^{1/3}$. This proves \eqref{eq:threeproducts}.
\end{proof}

The same proof applies to products of finitely many independent initial
tree moments whose leaves split at $p$, integrated against $\nu(\dd p)$.
The constant then depends on the total number of their initial leaves and on the
number of differentiations. Indeed, the initial mark measure is still
$\nu(\dd p)$, the number of factors to which the product rule applies
is at most the total number of initial leaves plus the number of differentiations,
and Steps~2 and~3 are unchanged.

\paragraph{The cavity covariance.}
In particular, \cref{lem:cav-threeproducts} applies to the covariance
between the squared conditional magnetization of the last spin and the
rest self-overlap, and to its Taylor remainder,
\[
 \mathcal C_{pp}(u)=\Cov_u\bigl(m_N(p)^2,A_p^-\bigr),\qquad
 \mathcal R_{\rm cav}=\int\nu(\dd p)\int_0^1(1-u)
                              \mathcal C_{pp}''(u)\dd u .
\]
Representing conditional expectations by leaves, we can write
$\mathcal C_{pp}=\Psi_p-\Psi^\epsilon_p\Psi^\tau_p$, where
\[
 \Psi_p(u)=\E_{4,p,u}[\epsilon_1\epsilon_2(R^-_{34}-p)],\qquad
 \Psi^\epsilon_p(u)=\E_{2,p,u}\epsilon_1\epsilon_2,\qquad
 \Psi^\tau_p(u)=\E_{2,p,u}(R^-_{12}-p).
\]
Indeed, by conditional independence given $\calG_p$, we have
$\Psi_p=\E_u[m_N(p)^2(A_p^--p)]$, $\Psi^\epsilon_p=\E_um_N(p)^2$, and
$\Psi^\tau_p=\E_u(A_p^--p)$. The constant $p$ does not change the
covariance. Note that the observables in these tree moments, including the center
$p$, do not depend on $u$. They are bounded by $2$, and the
trees have at most four leaves. By \cref{lem:streaming}, each of the
three moments belongs to $C^2([0,1])$, with first and second
derivatives bounded by $C\alpha$ and $C\alpha^2$. Then
$\mathcal C_{pp}\in C^2([0,1])$ for every deterministic profile, and
$|\mathcal C_{pp}''|$ is bounded by a constant depending only on
$\beta$, uniformly in $p$. In particular, $\mathcal R_{\rm cav}$ is well
defined, and by Taylor's formula with integral remainder, we have
\[
 \mathcal R_{\rm cav}=\int\bigl[\mathcal C_{pp}(1)-\mathcal C_{pp}(0)
   -\mathcal C_{pp}'(0)\bigr]\,\nu(\dd p).
\]

\begin{lemma}[Taylor remainder]\label{lem:remainder}
For every deterministic profile and every $u\in[0,1]$,
\begin{equation}\label{eq:C2derivative}
 \int|\mathcal C_{pp}''(u)|\,\nu(\dd p)\le C_\beta\mathcal M_3(u),
 \qquad
 |\mathcal R_{\rm cav}|\le C_\beta\sup_{0\le w\le1}\mathcal M_3(w).
\end{equation}
\end{lemma}

\begin{proof}
We have
\[
 \mathcal C_{pp}''=\Psi_p''-(\Psi^\epsilon_p)''\Psi^\tau_p
 -2(\Psi^\epsilon_p)'(\Psi^\tau_p)'-\Psi^\epsilon_p(\Psi^\tau_p)''.
\]
Each application of \eqref{eq:streaming} adds a factor
$\alpha/2\le\beta^2/2$ and one overlap defect, either inside a factor $D$
or in the independent tree of its last term. The spin factors have
absolute value one. The moments $\Psi_p$ and $\Psi^\tau_p$ contain one
defect, and $\Psi^\epsilon_p$ contains none. We count the defects before
differentiation and those added by the two differentiations.
\begin{center}
\small
\begin{tabular}{@{}lccc@{}}
\toprule
term & defects before & added & total\\
\midrule
$\Psi_p''$ & $1$ & $2$ & $3$\\
$(\Psi^\epsilon_p)''\Psi^\tau_p$ & $0+1$ & $2+0$ & $3$\\
$(\Psi^\epsilon_p)'(\Psi^\tau_p)'$ & $0+1$ & $1+1$ & $3$\\
$\Psi^\epsilon_p(\Psi^\tau_p)''$ & $0+1$ & $0+2$ & $3$\\
\bottomrule
\end{tabular}
\end{center}
Products of moments, and the last term of \eqref{eq:streaming}, are
products over independent trees. Every term of $\mathcal C_{pp}''$ is
therefore $(\alpha/2)^2$ times a term of the kind treated in
\cref{lem:cav-threeproducts}, or in its extension to products of
independent initial moments stated after its proof, with at most four
initial leaves and $j=2$. It is multiplied by a bounded multiplicity and
by factors of absolute value at most one. Integrating $p$ against $\nu$
and using \cref{lem:cav-threeproducts} and this extension, we obtain the first bound in
\eqref{eq:C2derivative} with $C_\beta=C\beta^4$. We integrate it against
$(1-u)\dd u$ and use Tonelli's theorem. This gives
$|\mathcal R_{\rm cav}|\le C\beta^4\int_0^1(1-u)\mathcal M_3(u)\dd u
\le C\beta^4\sup_w\mathcal M_3(w)$, which is the second bound.
\end{proof}
